\chapter*{برای مطالعه‌ی بیشتر}
\addcontentsline{toc}{chapter}{برای مطالعه‌ی بیشتر}
\markboth{برای مطالعه‌ی بیشتر}{برای مطالعه‌ی بیشتر}

بخش منابع تمام آثار استناد‌شده در متن را فهرست می‌کند. این بخش فراتر از آن‌ها
می‌رود --- به کتاب‌ها و مقالاتی اشاره می‌کند که اگر ایده‌های این کتابچه
سوال‌هایی را در ذهنت باز کرده که می‌خواهی دنبالشان کنی، ارزش خواندن دارند.

\medskip

\textbf{درباره‌ی داده، اطلاعات و دانش.}
کتاب
\begin{latin}\textit{Information: A Very Short Introduction}\end{latin}
نوشته‌ی لوچیانو فلوریدی (آکسفورد، ۲۰۱۰) دقیق‌ترین و در عین حال
قابل‌دسترس‌ترین تحلیل از مفهوم اطلاعات است.
مقاله‌ی کوتاه راسل اکاف با عنوان
\begin{latin}``From Data to Wisdom'' (1989)\end{latin}
که در
\begin{latin}\textit{Journal of Applied Systems Analysis}\end{latin}
منتشر شده، هرم \lr{DIKW} را معرفی کرده و تا امروز نقطه‌ی شروع استاندارد
بحث درباره‌ی رابطه‌ی داده، اطلاعات، دانش و خرد است.

\medskip

\textbf{درباره‌ی بازنمایی و نشانه‌شناسی.}
مدخل «بازنمایی ذهنی» در دایره‌المعارف فلسفه‌ی استنفورد
(\url{https://plato.stanford.edu/entries/mental-representation/})
تحلیل فلسفی دقیقی از مفهوم بازنمایی ارائه می‌دهد.
برای نشانه‌شناسی به معنای گسترده‌تر، کتاب
\begin{latin}\textit{A Theory of Semiotics}\end{latin}
نوشته‌ی اومبرتو اکو (دانشگاه ایندیانا، ۱۹۷۶) متن کلاسیک این حوزه است.

\medskip

\textbf{درباره‌ی بستر فیزیکی.}
جلدهای بعدی مجموعه‌ی \textit{آرلیز: صفر تا بیت} دقیقاً سوالی را که
فصل سوم این کتابچه با آن تمام می‌شود دنبال می‌کنند: اینکه چطور یک
بازنمایی انتزاعی به یک سیگنال فیزیکی تبدیل می‌شود. این جلدها قدم بعدی
پیشنهادی هستند و در مخزن پروژه در دسترسند:
\url{https://github.com/papyrxis/arliz}